\begin{center}
  {\heiti\zihao{2}\bfseries \PaperTitle}
\end{center}
\vspace{0.5em}

\section*{摘要}

本文针对乙醇催化偶合制备 $\mathrm{C}_4$ 烯烃过程中催化剂组合、温度及反应时间对反应性能的影响问题，基于实验数据的分组结构建立分问题模型。针对不同任务分别采用温度响应分析、因素匹配比较、收率优化和追加实验设计方法，避免将非正交小样本实验简单视为同质独立样本。

针对\textbf{问题一}，通过单催化剂组合内温度趋势分析和附件2时间序列稳定性分析，建立低复杂度经验响应模型，刻画乙醇转化率、$\mathrm{C}_4$ 烯烃选择性及稳定性变化规律。

针对\textbf{问题二}，在不把温度重复记录视为独立催化剂设计的前提下，建立基于温度分层匹配与跨背景效应对比的催化因素影响分析。21 个组合在 250--350 ℃端点上的乙醇转化率与 $\mathrm{C}_4$ 烯烃选择性均呈正增量，平均分别提高 24.46、15.30 个百分点；严格匹配和局部四格分析进一步揭示了催化因素影响的背景依赖性。

针对\textbf{问题三}，通过构造 $\mathrm{C}_4$ 烯烃收率指标，在实验约束范围内寻找高收益条件，并结合模型稳定性分析评价优化结果可靠性。

本文模型充分利用实验设计中的控制变量结构，具有较好的可解释性和可追溯性；但由于原实验规模有限，所得因素影响主要适用于已有实验范围。

\noindent\textbf{关键词：} 乙醇偶合；催化剂组合；温度响应；严格匹配分析；实验优化
